\documentclass{ctexart}
\usepackage{../../Styles/mystyle-report}

\setlist[enumerate,1]{label=(\arabic*),itemsep=4pt,leftmargin=*}
\setlist[enumerate,2]{label=(\roman*),itemsep=3pt,leftmargin=2em}

% 去掉超链接的红框
\hypersetup{hidelinks}

\begin{document}

% 生成目录
\tableofcontents
\newpage

\section{计量经济学的作用？/计量经济模型有哪些应用？}
\begin{enumerate}
\item 结构分析；
\item 经济预测；
\item 政策评价；
\item 检验与发展经济理论。
\end{enumerate}

\section{确定解释变量的注意点}
\begin{enumerate}
\item 符合经济理论;

\item 数据可得可量化;

\item 避免高度共线性;

\item 保证外生性;

\item 精简适度。
\end{enumerate}

\section{建立计量模型的四个阶段？}
\begin{enumerate}
\item 设定模型形式；
\item 估计参数；
\item 检验模型；
\item 模型的应用。
\end{enumerate}

\section{随机误差项产生的原因？}
\begin{enumerate}
\item 省略变量；
\item 随机因素；
\item 测量误差；
\item 模型设定偏差；
\item 数据归并误差。
\end{enumerate}

\section{古典线性回归模型的基本假定是什么？/简述最小二乘法的基本假定}
\begin{enumerate}
\item 零均值假定：随机误差均值为0，即 $E(u_i)=0$；
\item 同方差假定：误差项方差恒定，即 $\text{Var}(u_i)=\sigma^2$；
\item 无自相关假定：不同样本的误差项无相关性，即 $\text{Cov}(u_i,u_j)=0\ (i\neq j)$；
\item 外生性假定：解释变量与误差项不相关，即 $\text{Cov}(X_i,u_i)=0$；
\item 正态性假定：误差项服从正态分布，即 $u_i \sim N(0,\sigma^2)$。
\end{enumerate}

\section{最小二乘估计原理？}
最小二乘法核心是最小化残差平方和，即使
$$
Q = \sum_{i=1}^n e_i^2 = \sum_{i=1}^n (Y_i - \hat{Y}_i)^2 
$$
达到最小，通过求偏导为0求解得到模型参数的估计量。

\section{t检验的步骤？}
\begin{enumerate}
\item 提出假设：原假设 $H_0: \beta_i=0$，备择假设 $H_1: \beta_i\neq0$；
\item 构造统计量：$t = \frac{\hat{\beta}_i}{S(\hat{\beta}_i)},$
当$H_0$ 成立时，$t \sim t(n-k-1)$,其中$k$为模型中解释变量个数；
\item 确定临界值：给定显著性水平 $\alpha$，查表得到临界值 $t_{\frac{\alpha}{2}}(n-k-1)$；
\item 判断显著性：若 $|t|>t_{\frac{\alpha}{2}}(n-k-1)$，拒绝 $H_0$，变量影响显著；反之则影响不显著。
\end{enumerate}

\section{回归分析的主要内容是什么？}
\begin{enumerate}
\item 利用样本数据估计模型未知参数，得到样本回归方程；
\item 检验回归方程整体与单个参数的统计显著性，验证模型可靠性；
\item 利用回归方程对被解释变量进行预测或结构分析。
\end{enumerate}

\section{说明利用样本判定系数 $R^2$，为什么能判断回归直线与样本观察值的拟合优度？}
判定系数 $R^2$ 代表模型能解释的变异占总变异的比例：
\begin{enumerate}
\item $R^2$ 越接近1，模型解释的变异占比越高，回归直线对样本的拟合程度越好；
\item $R^2$ 越接近0，模型解释能力越弱，拟合程度越差。
\end{enumerate}

\section{多元回归模型基本假定}
\begin{enumerate}
\item 零均值假定：随机误差均值为0，即 $E(u_i)=0$；
\item 同方差假定：误差项方差恒定，即 $\text{Var}(u_i)=\sigma^2$；
\item 无自相关假定：不同样本的误差项无相关性，即 $\text{Cov}(u_i,u_j)=0\ (i\neq j)$；
\item 外生性假定：解释变量与误差项不相关，即 $\text{Cov}(X_i,u_i)=0$；
\item 正态性假定：误差项服从正态分布，即 $u_i \sim N(0,\sigma^2)$；
\item 解释变量之间不存在完全多重共线性。
\end{enumerate}

\section{多元回归分析中为何要使用调整的判定系数}
为避免因解释变量个数增加而导致判定系数虚增的假象，故采用调整的判定系数，通过对自由度进行惩罚来剔除变量个数对拟合优度的影响。

\section{对于多元线性回归模型，为什么在进行了总体显著性F检验之后，还要对每个偏回归系数进行是否为0的t检验？}
因为F检验只能证明模型中所有解释变量对被解释变量有整体显著性，但无法指明在控制其他变量后哪些变量仍有独立的显著影响，所以必须进一步对每个偏回归系数进行t检验。

\section{简述多元回归模型的整体显著性检验（F检验）的步骤}
\begin{enumerate}
\item 提出假设:
$$
\text{原假设}H_0: \beta_1=\beta_2=\dots=\beta_k=0, \quad \text{备择假设}H_1: \text{至少有一个 } \beta_i \neq 0,
$$
其中 $\beta_i$ 为各解释变量的偏回归系数。

\item 构造统计量：
$$
F = \frac{ESS/k}{RSS/(n-k-1)}
$$
其中 $ESS$ 为回归平方和，$RSS$ 为残差平方和；$k$ 为解释变量个数，$n$ 为样本容量。当$H_0$ 成立时，$F \sim F(k,\ n-k-1)$。

\item 确定临界值：给定显著性水平 $\alpha$，查表得到临界值 $F_\alpha(k,\ n-k-1)$；

\item 判断显著性：若 $F > F_\alpha$，拒绝 $H_0$，回归方程整体线性关系显著；反之则整体不显著。
\end{enumerate}

\section{举例说明经济现象中的异方差性}

居民收入与消费支出。高收入群体消费选择差异大，低收入群体消费相对集中，导致随收入增加，消费支出的波动增大，即存在异方差性。

\section{说明加权最小二乘法的基本思想}

通过对原模型变换，对方差小的观测值赋予较大权重，方差大的观测值赋予较小权重，权重通常取为误差项方差估计值的倒数。使加权后的模型满足同方差假定，再使用普通最小二乘法估计参数。

\section{说明戈德菲尔德——匡特（Goldfeld-Quandt）检验的步骤}
\begin{enumerate}
\item 设原假设 $H_0: \sigma_1^2 = \sigma_2^2$（随机误差项具有同方差性），

备择假设 $H_1: \sigma_1^2 \neq \sigma_2^2$（随机误差项具有异方差性）；

\item 将样本按某一解释变量升序排列；
\item 去掉中间$c$个观测值，将剩余观测值分为两个子样本；
\item 对两个子样本分别进行OLS回归，计算各自的残差平方和，将其中较小者记为$RSS_1$，较大者记为$RSS_2$；
\item 构造F统计量：$F = \dfrac{RSS_2 / (n_2 - k-1)}{RSS_1 / (n_1 - k-1)}$,其中$k$为模型中解释变量个数,$n_1,n_2$分别为两个子样本容量。当$H_0$成立时,$F\sim F(n_2-k-1,n_1-k-1)$；
\item 给定显著性水平$\alpha$，若$F$大于临界值$F_{\alpha}(n_2-k-1,n_1-k-1)$，则拒绝原假设，认为存在异方差。反之则不存在异方差。
\end{enumerate}

\section{戈里瑟（Glejser）检验的步骤}
\begin{enumerate}
\item 用普通最小二乘法估计原模型，得到残差绝对值$|\hat{e}_i|$；
\item 用$|\hat{e}_i|$分别对某个解释变量$X_i$的多种函数形式(如$\sqrt{X_i},\frac{1}{X_i},X_i^2$等)进行辅助回归；
\item 通过辅助回归系数的显著性判断异方差性；若某种函数形式的系数在$t$检验下显著不为零，则说明模型存在异方差。
\end{enumerate}

\section{杜宾—瓦森（D-W）检验适用条件}
\begin{enumerate}
\item 模型包含截距项；
\item 扰动项为一阶自相关形式；
\item 解释变量非随机，且模型不含滞后被解释变量；
\item 样本为连续时间序列，无缺失数据。
\end{enumerate}

\section{杜宾—瓦森（D-W）检验步骤/D-W的具体内容}

DW 检验是检验随机误差项一阶自相关的常用方法，具体步骤如下：
\begin{enumerate}
\item 设原假设 $H_0: \rho = 0$（无一阶自相关），备择假设 $H_1: \rho \neq 0$；
\item 用最小二乘法估计模型，得到残差序列 $e_t$；
\item 计算DW统计量：$DW = \frac{\sum\limits_{t=2}^{n}(e_t - e_{t-1})^2}{\sum\limits_{t=1}^{n} e_t^2}$；
\item 给定显著性水平，查临界值 $d_L$、$d_U$；
\item 对比DW与临界区间，判断是否存在自相关。
\end{enumerate}

\section{杜宾—瓦森（D-W）检验的缺点}
\begin{enumerate}
\item 仅能检验一阶自相关，无法检验高阶自相关及其他形式的自相关；

\item 当DW值在临界值\(d_L\)与\(d_U\)之间时,无法判定自相关性；

\item 模型中若包含滞后被解释变量，检验失效；

\item 要求样本为连续时间序列，数据不能存在缺失。
\end{enumerate}

\section{自相关来源}
\begin{enumerate}
\item 经济变量的惯性与趋势性；
\item 模型设定偏误：遗漏变量、函数形式错误；
\item 数据加工（插值、平滑）导致的伪自相关；
\item 经济行为的滞后效应。
\end{enumerate}

\section{多重共线性的定义与识别方法}

多元线性回归模型中，两个及以上解释变量之间存在精确的或近似的线性相关关系，即为多重共线性。识别多重共线性的方法有：
\begin{enumerate}
\item 简单相关系数法;

\item 方差膨胀因子（VIF）法;

\item 经验判断法;

\item 逐步回归法.
\end{enumerate}

\section{多重共线性的原因}
\begin{enumerate}
\item 经济变量之间的内在联系；
\item 经济变量存在着同方向变动的趋势；
\item 样本资料匮乏；
\item 模型中引入了重复或多余的变量；
\item 将某些解释变量的滞后值作为解释变量引入了模型之中。
\end{enumerate}

\section{多重共线性的影响后果}
\begin{enumerate}
\item OLS估计量仍无偏，但方差大幅膨胀，估计精度下降；
\item t统计量变小，变量显著性检验失效，易出现整体显著、单个不显著；
\item 参数估计值不稳定，样本微小变动会造成系数大幅波动，符号易违背经济理论；
\item 难以区分单个解释变量对被解释变量的独立作用；
\item 若预测期共线性关系不变，则模型仍可用于预测。
\end{enumerate}

\section{完全多重共线性对OLS估计量的影响有哪些？}
\begin{enumerate}
\item 矩阵$X'X$奇异不可逆，OLS的参数估计量$\hat{\beta} = (X'X)^{-1}X'Y$不存在；
\item 参数方差无穷大，完全无法识别各变量独立影响；
\item 无法开展t检验、F检验，无法确定模型估计结果。
\end{enumerate}


\section{克服多重共线性的方法}
\begin{enumerate}
\item 剔除次要变量;

\item 扩大样本容量;

\item 变换模型形式;

\item 施加先验约束;

\item 逐步回归筛选变量。
\end{enumerate}

\section{对数线性模型的优点}
\begin{enumerate}
\item 系数直接表示弹性，经济含义清晰直观；

\item 可将幂函数形式的非线性关系转化为线性形式，便于使用普通最小二乘法估计；

\item  对变量取对数可压缩数据量级，一定程度上缓解异方差问题；

\item  能部分降低解释变量之间的多重共线性程度。
\end{enumerate}



\end{document}